% =====================================================================
% I. INTRODUCCIÓN
% =====================================================================
\section{Introducci\'on}

La digitalizaci\'on del patrimonio musical constituye una necesidad
creciente para archivos, bibliotecas, academias y agrupaciones musicales
que conservan partituras f\'isicas en papel. El reconocimiento \'optico de
m\'usica (OMR, por sus siglas en ingl\'es) es la disciplina que aborda la
conversi\'on autom\'atica de im\'agenes de partituras en representaciones
simb\'olicas editables, y ha sido objeto de investigaci\'on durante m\'as de
cinco d\'ecadas \cite{omr2020}, evolucionando desde sistemas basados en
reglas hasta \emph{pipelines} de aprendizaje profundo que transcriben
partituras completas de forma autom\'atica
\cite{rebelo2012}\cite{calvo2018}\cite{riosvila2026}.

A pesar de estos avances, los sistemas OMR modernos cometen errores que
impiden su uso directo en entornos reales: notas con alturas o duraciones
incorrectas, compases desbalanceados, armaduras inconsistentes y s\'imbolos
omitidos \cite{calvo2018}. Los modelos de reconocimiento se estudian
predominantemente desde la precisi\'on de su salida autom\'atica, mientras
que la integraci\'on del reconocimiento en el trabajo pr\'actico de un
transcriptor ha recibido mucha menos atenci\'on \cite{rizo2023}. Como
consecuencia, quien digitaliza partituras termina revisando manualmente
cada s\'imbolo en un editor musical convencional, un proceso tan lento como
la transcripci\'on manual, o acepta resultados err\'oneos sin verificarlos.
En el caso de los manuscritos musicales modernos, esta situaci\'on se
agudiza, pues la variabilidad de la escritura degrada a\'un m\'as la
confiabilidad del reconocimiento autom\'atico \cite{ricordi2024}. Del mismo
modo, el aprovechamiento de las correcciones del usuario para mejorar el
sistema mediante ciclos de aprendizaje activo sigue siendo poco explorado
en OMR, donde los primeros estudios arrojan resultados incluso
contradictorios \cite{mlsp2025}.

En el contexto global, instituciones conservan colecciones de partituras
impresas y escritas a mano que requieren digitalizaci\'on para su
preservaci\'on, consulta y difusi\'on. La transcripci\'on manual de estas
colecciones representa un costo de tiempo considerable, y las herramientas
de reconocimiento existentes no ofrecen un flujo de trabajo que garantice
la calidad del resultado. La digitalizaci\'on de colecciones musicales ha
sido abordada recientemente desde la arch\'istica con metodolog\'ias de
anotaci\'on y clasificaci\'on autom\'atica \cite{ricordi2024}; a\'un as\'i,
faltan herramientas de transcripci\'on asistida pensadas para el trabajo
del transcriptor.

Para la evaluaci\'on del sistema se dispone de conjuntos de datos
p\'ublicos de referencia que cubren el alcance del proyecto. PrIMuS y
Camera-PrIMuS \cite{calvo2018} proveen 87\,678 incipits de partituras
impresas monof\'onicas, el segundo con distorsiones que simulan fotograf\'ias
tomadas en condiciones reales, y \emph{ground truth} simb\'olico en
m\'ultiples formatos. El Sheet Music Benchmark (SMB)
\cite{smb2026} estandariza la evaluaci\'on con 685 p\'aginas que incluyen
monofon\'ia y \emph{pianoform}, divisiones predefinidas y una m\'etrica de
error detallada por elemento musical (OMR-NED). Para el subconjunto
manuscrito, MUSCIMA++ \cite{muscima2017} ofrece m\'as de 91\,000 s\'imbolos
anotados sobre partituras escritas a mano por m\'usicos actuales. La
disponibilidad de \emph{ground truth} anotado en estos conjuntos permite
evaluar la plataforma de forma cuantitativa y reproducible, siguiendo los
protocolos de la literatura de OMR. De manera complementaria, se contempla
un piloto sobre un conjunto reducido de partituras reales (fotograf\'ias)
para evaluar la usabilidad del flujo completo de correcci\'on, sin requerir
\emph{ground truth} anotado.

Desde el punto de vista disciplinar, un sistema de digitalizaci\'on
asistida de partituras puede estructurarse sobre tres pilares: el
\emph{pipeline} cl\'asico de OMR (preprocesado, detecci\'on de pentagramas,
detecci\'on y clasificaci\'on de s\'imbolos, ensamblaje y reconstrucci\'on
musical \cite{omr2020}), la validaci\'on de la salida
mediante reglas de teor\'ia musical \cite{cuthbert2010}, y la correcci\'on
asistida por un humano en el ciclo \cite{amershi2014}. La validaci\'on
autom\'atica permite detectar errores estructurales (m\'etrica, duraciones,
armaduras) sin intervenci\'on humana, reduciendo el espacio de revisi\'on a
los puntos sospechosos. El aprendizaje activo complementa el ciclo al
seleccionar qu\'e correcciones del usuario tienen mayor valor para mejorar
el reconocimiento \cite{settles2009}.

En este contexto, el presente proyecto propone el dise\~no e
implementaci\'on de \textbf{Cadenza}, una plataforma de digitalizaci\'on
asistida de partituras que integra cuatro componentes: un modelo OMR base
(oemer \cite{oemer2021}) encargado de la transcripci\'on autom\'atica a
MusicXML; un motor de validaci\'on basado en reglas musicales que detecta y
se\~nala errores probables; una interfaz de correcci\'on asistida que
permite al usuario revisar y corregir la transcripci\'on sobre la imagen
original, con reproducci\'on autom\'atica de la partitura digitalizada; y un ciclo
de aprendizaje activo que aprovecha las correcciones
del usuario para priorizar la revisi\'on y mejorar progresivamente el
sistema. La plataforma se orienta a partituras impresas y escritas a mano
actuales, monof\'onicas o de piano simple, excluyendo orquestaciones
completas y manuscritos hist\'oricos, cuyo tratamiento excede el alcance de
este trabajo.
